\scanpage{45}
\mspage{- 23 -}
% block: top; line 1; display
\[
Du_i(x;z_1,\mathinner{\ldotp\ldotp},z_q)=f_i(x;u_1(x;z_1,\mathinner{\ldotp\ldotp},z_q),\mathinner{\ldotp\ldotp},u_p(x;z_1,\mathinner{\ldotp\ldotp},z_q);z_1,\mathinner{\ldotp\ldotp},z_q)
\tag*{(12)}
\]
% block: top; line 2
(D \uline{dérivation par rapport à} x, $1\handwritten{\leq}i\handwritten{\leq}p$) \uline{et}
% block: top; line 3; display
\[
u_i(a;z_1,\mathinner{\ldotp\ldotp},z_q)=b_i\qquad(1\handwritten{\leq}i\handwritten{\leq}p)\quad\text{\uline{pour tout}}\quad(z_1,\mathinner{\ldotp\ldotp},z_q)\handwritten{\in}W
\tag*{(13)}
\]
% block: top; lines 4-7
Par translation, on peut se ramener au cas où $a=b_j=c_k=0$ pour $1\handwritten{\leq}j\handwritten{\leq}p$ et $1\handwritten{\leq}k\handwritten{\leq}q$ ; chacune des fonctions $f_i$ est alors égale à la somme d'une série entière sommable dans un polycylindre ouvert $B\handwritten{\times}P\handwritten{\times}Q$, où $B\handwritten{\subset}K$, $P\handwritten{\subset}K^p$ et $Q\handwritten{\subset}K^q$
% block: middle; lines 8-9; display
\[
\begin{aligned}
f_i(x;y_1,\mathinner{\ldotp\ldotp},y_p;z_1,\mathinner{\ldotp\ldotp},z_q)&=\handwritten{\sum}g_i(n;r_1,\mathinner{\ldotp\ldotp},r_p;\handwritten{\unsure{\lambda|\Lambda}}_1,\mathinner{\ldotp\ldotp},\handwritten{\lambda}_q)\,\ldotp\\
&\quad\ldotp x^ny_1^{r_1}\mathinner{\ldotp\ldotp}y_p^{r_p}z_1^{\handwritten{\lambda}_1}z_2^{\handwritten{\lambda}_2}\ldots z_q^{\handwritten{\lambda}_q}.
\end{aligned}
\tag*{(14)}
\]
% block: middle; lines 10-11
Pour établir l'existence des fonctions $u_i$, nous devons prouver qu'il existe p séries entières sans terme indépendant de x
% block: middle; line 12; display
\[
u_i(x;z_1,\mathinner{\ldotp\ldotp},z_q)=\handwritten{\sum}h_i(m;\handwritten{\mu}_1,\mathinner{\ldotp\ldotp},\handwritten{\mu}_q)x^mz_1^{\handwritten{\mu}_1}\ldots z_q^{\handwritten{\mu}_q}
\]
% block: middle-bottom; lines 13-17
sommables dans un polycylindre $U\handwritten{\times}W$ de $K^{q+1}$, substituables dans les séries entières (14) et satisfaisant à (12). Si le problème est possible, les séries entières par rapport à x et aux $z_{\handwritten{\ell}}$ $(1\handwritten{\leq\ell\leq}q)$ obtenues par cette substitution dans les deux membres de (12) doivent avoir mêmes coefficients, ce qui donne les conditions
% block: bottom; lines 18-20; display
\[
\left.\begin{aligned}
h_i(0;\handwritten{\mu}_1,\mathinner{\ldotp\ldotp},\handwritten{\mu}_q)&=0\\
mh_i(m;\handwritten{\mu}_1,\ldots,\handwritten{\mu}_q)&=\handwritten{\sum_{n,(r_j),(\lambda_\ell)}}\unsure{G|g}_i(n;r_1,\mathinner{\ldotp\ldotp},r_p;\handwritten{\lambda}_1,\ldots,\handwritten{\lambda}_q)\,\ldotp\\
&\quad\ldotp\handwritten{\prod_{j=1}^{p}\prod_{k=1}^{r_j}}h_i(s_{kj};\handwritten{\nu}_{kj1},\ldots,\handwritten{\nu}_{kjq})
\end{aligned}\handwritten{\right\}}
\tag*{(15)}
\]
% block: bottom; lines 21-24
pour tout entier $m\handwritten{>}0$ et tout système d'entiers $\handwritten{\mu_\ell\geq}0$ $(1\handwritten{\leq\ell\leq}q)$ ; la sommation du second membre des relations (15) doit être étendue à tous les systèmes d'entiers $\handwritten{\geq}0$, $n,r_j,\handwritten{\lambda_\ell},s_{kj}$ et $\handwritten{\nu}_{kj\handwritten{\ell}}$ satisfaisant aux conditions
% archive sheet number: 44
